\section{作者合作网络分析}

\begin{reportbox}{分析目标}
作者合作网络用于识别核心研究者、跨领域连接者和合作团簇。它是知识图谱查询和论文推荐解释的重要依据, 也是数据分析报告中最能体现“科研智能体”差异化的图分析部分。
\end{reportbox}

\subsection{网络构建}

\begin{itemize}
  \item \textbf{节点}: 作者, 后续可扩展为机构、国家、领域、关键词。
  \item \textbf{边}: 两位作者共同出现在同一篇论文中。
  \item \textbf{边权}: 合作次数、近年合作次数、代表论文质量。
  \item \textbf{节点属性}: 论文数、领域分布、活跃年份、中心性、社区编号。
\end{itemize}

\subsection{核心指标}

\begin{center}
\begin{tabularx}{\textwidth}{>{\bfseries}p{32mm}X X}
\toprule
指标 & 含义 & 报告解释方式 \\
\midrule
Degree & 作者直接合作人数 & 识别高合作度作者 \\
Betweenness & 作者连接不同团簇的桥梁能力 & 识别跨领域连接者 \\
Community & 合作社区编号 & 发现研究团队与方向群落 \\
Temporal Activity & 作者按年份的活跃情况 & 判断作者是否仍处于活跃期 \\
\bottomrule
\end{tabularx}
\end{center}

\subsection{合作社区与高频合作关系}

\begin{figure}[htbp]
  \centering
  \figasset[0.9\textwidth]{author_network_scale_funnel.png}
  \caption{2022--2026 作者合作网络规模漏斗。}
\end{figure}

\begin{figure}[htbp]
  \centering
  \figasset[0.88\textwidth]{author_collaboration_network.png}
  \caption{2022--2026 近五年作者合作社区气泡图。}
\end{figure}

\begin{figure}[htbp]
  \centering
  \figasset[0.9\textwidth]{top_author_collaborations.png}
  \caption{2022--2026 近五年高频作者合作关系。}
\end{figure}

作者合作边由论文作者列表两两组合得到。为避免一篇大团队论文生成过多虚高边, 静态报告图仅纳入作者数不超过 12 的论文, 并对每篇论文内的合作关系使用分数化权重。当前作者社区图使用 25,000 条高权重合作边构建核心图, 既保留 50 万级候选边的规模依据, 又避免 PDF 静态图退化成不可读的全量网络。

\clearpage
